% 5 分钟汇报版：在 concise 原稿基础上突出重点、补足结果解释
% 编译：xelatex main.tex 两遍（TeXLive / XeLaTeX）
\PassOptionsToPackage{quiet}{fontspec}
\PassOptionsToPackage{dvipsnames}{xcolor}
\documentclass[10pt,aspectratio=169]{beamer}
\usepackage{colortbl}
\usepackage{tikz}
\usetikzlibrary{arrows.meta,positioning}
\usepackage{zju_beamer}
\graphicspath{{./figures/}}
\zjusectioncolor{blue}

\title[GenAI 课程项目报告]{面向 CODESYS 平台的 ST 代码自动生成系统}
\subtitle{——GenAI 课程大作业}
\author[姓名]{姓名} % TODO: 原文未给出作者信息，请补全
\institute[浙江大学]{浙江大学}
\date{\today}

% 大数字强调块
\newcommand{\statbox}[2]{%
  \begin{minipage}[t]{0.31\linewidth}\centering
    {\color{blue!65!black}\bfseries\Large #1}\par\vspace{1pt}
    {\scriptsize #2}%
  \end{minipage}}

\begin{document}

\begin{frame}
  \titlepage
\end{frame}

% ============================================================
\section{问题总览}

\begin{frame}
  \frametitle{问题总览}

  \begin{columns}[T]
    \begin{column}{0.48\textwidth}
      \begin{block}{目标}
        输入自然语言需求 + 接口表 + LTL 规约；
        输出可直接在 CODESYS 编译并运行的 \textbf{ST}
      \end{block}
      \begin{alertblock}{难点不在“像不像”}
        \begin{itemize}
          \item 语法与工程结构：DUT、Standard 库、调用图
          \item 行为语义：状态机、定时器、边沿、优先级
          \item 评测可信：必须\textbf{真实执行}
        \end{itemize}
      \end{alertblock}
    \end{column}
    \begin{column}{0.50\textwidth}
      \begin{block}{交付}
        数据与规约（30 题 / 75 条 LTL / 54 组用例）$\cdot$
        自研 Agent harness $\cdot$ CODESYS 编译 + 仿真服务 $\cdot$
        统一评测管线
      \end{block}
      \begin{block}{四种设置}
        \{Codex CLI, 自研\} $\times$ \{DeepSeek, Qwen\}，各 30 题，
        统一 5 轮修复预算，\texttt{tool\_error} 已清零
      \end{block}
    \end{column}
  \end{columns}

\end{frame}

% ============================================================
\section{数据与规约}

\begin{frame}
  \frametitle{数据集与质量基线}

  \begin{columns}[T]
    \begin{column}{0.46\textwidth}
      \centering
      \small
      \begin{tabular}{@{}ll@{}}
        \hline
        维度 & 数量 \\
        \hline
        总题数 & 30（公共 20 + 自建 10） \\
        难度 & easy 9 / medium 12 / hard 9 \\
        形式化规约 & 75 条（2.5 条/题） \\
        测试用例 & 54 组、199 个扫描步 \\
        \hline
      \end{tabular}

      \vspace{2mm}
      \raggedright\scriptsize
      每题包含：\texttt{FUNCTION\_BLOCK} 接口（含 DUT）+
      中英需求 + LTL + 行为用例 + 可编译参考实现
    \end{column}
    \begin{column}{0.50\textwidth}
      \begin{itemize}
        \item \textbf{自建题}：交通灯、水箱、自动门、分拣、主备泵、报警闭锁、
              星三角、停车场、恒压供水、电梯——均为真实 PLC 场景
        \item \textbf{多样性}：离散制造/流程工业/电机/安全/队列/字符串；
              含 DUT 3 题、\texttt{VAR\_IN\_OUT} 4 题
        \item \textbf{规约}：互斥、响应、保持、范围不变量、错误码
        \item \textbf{用例}：正常流程、边界、优先级冲突、错误处理；
              \textbf{每步写明测试目的}并关联到 LTL 规约
        \item \textbf{自检}：参考实现 \textbf{30/30} 编译且 \textbf{30/30}
              通过自己的用例（首轮曾暴露 16 处用例缺陷）
      \end{itemize}
    \end{column}
  \end{columns}

\end{frame}

% ============================================================
\section{Agent 框架}

\begin{frame}
  \frametitle{架构：单 Agent 反馈闭环}

  \begin{columns}[T]
    \begin{column}{0.45\textwidth}
      \centering
      \begin{tikzpicture}[
          x=1cm, y=0.88cm,
          box/.style={draw=structure!80, fill=structure!8, rounded corners=2pt,
            align=center, font=\scriptsize, text width=3.5cm, inner sep=1.6pt},
          fin/.style={draw=green!45!black, fill=green!10, rounded corners=2pt,
            align=center, font=\scriptsize, text width=3.5cm, inner sep=1.6pt},
          a/.style={-{Latex[length=1.3mm]}, semithick},
          fb/.style={semithick, draw=red!60!black}]
        \node[box]  (t)  at (0,0)      {任务：需求 + 接口 + LTL};
        \node[box]  (p)  at (0,-1.05)  {提示构造 + 模型生成};
        \node[box]  (v)  at (0,-2.10)  {解析 + 结构校验};
        \node[box]  (k)  at (0,-3.15)  {组装工程 + CODESYS 编译};
        \node[box]  (s)  at (0,-4.20)  {仿真行为测试};
        \node[fin]  (ok) at (0,-5.25)  {成功};
        \foreach \u/\d in {t/p, p/v, v/k, k/s, s/ok}
          \draw[a] (\u) -- (\d);
        \draw[fb] (v.west) -- (-2.45,-2.10);
        \draw[fb] (k.west) -- (-2.45,-3.15);
        \draw[fb] (s.west) -- (-2.45,-4.20);
        \draw[fb] (-2.45,-4.20) -- (-2.45,-1.05);
        \draw[-{Latex[length=1.3mm]}, semithick, draw=red!60!black]
          (-2.45,-1.05) -- (p.west);
        \node[font=\tiny, text=red!60!black, rotate=90, fill=white, inner sep=1pt]
          at (-2.45,-3.15) {结构化反馈};
      \end{tikzpicture}
    \end{column}
    \begin{column}{0.52\textwidth}
      \begin{itemize}
        \item 一次生成完整工程，每轮整体重写，不做局部补丁
        \item \textbf{结构校验前置}：格式/接口错误不消耗 CODESYS 冷启动
        \item \textbf{上下文}：仅保留最近 2 轮；错误按 (code, text) 去重
        \item \textbf{终止}：通过 / 5 轮 / 代码无变化 / 失败签名重复 / 工具故障
        \item 测试用例\textbf{对模型不可见}，仿真基于 CODESYS 官方运行时
        \item 两道闸门：接口校验零成本，编译约 13 s，仿真约 24 s
      \end{itemize}
    \end{column}
  \end{columns}

\end{frame}

% ============================================================
\section{结果与解释}

\begin{frame}
  \frametitle{实验设置与公平性}

  \centering
  \begin{tabular}{lcc}
    \hline
     & deepseek-v4.1-flash-1017 & qwen3.7-flash \\
    \hline
    现有框架（Codex CLI） & 实验 1 & 实验 2 \\
    自研 harness          & 实验 3 & 实验 4 \\
    \hline
  \end{tabular}

  \vspace{2.5mm}
  \begin{itemize}
    \item Codex CLI 包装为与自研相同的 \texttt{chat(messages)} 客户端：
          \textbf{提示、反馈文本、轮次上限、编译测试管线完全一致}
    \item 差异只在生成组织方式（Codex 自带 agent 循环与工具上下文），
          因此结果差异只能归因于框架本身
    \item 同一 DMXAPI 网关、同一模型快照、同一 CODESYS 模板与 worker
    \item \textbf{指标}：首次编译、最终编译、测试通过率、token、墙钟时间、轮次
    \item Codex 的工具调用开销计入 token 与耗时——这是“现有框架”的真实成本
  \end{itemize}

\end{frame}

\begin{frame}
  \frametitle{总体结果：DeepSeek 打平，Qwen 拉开差距}

  \centering
  \small
  \setlength{\tabcolsep}{6pt}
  \renewcommand{\arraystretch}{1.25}
  \begin{tabular}{lccccc}
    \hline
    设置 & 首次编译 & 最终编译 & 测试通过 & 平均轮次 & 平均耗时 \\
    \hline
    自研 + DeepSeek  & 96.7\%  & 100.0\% & \textbf{83.3\%} & 1.77 & 163 s \\
    自研 + Qwen      & 76.7\%  & 93.3\%  & \textbf{66.7\%} & 2.53 & 295 s \\
    Codex + DeepSeek & 100.0\% & 100.0\% & \textbf{83.3\%} & 1.70 & 342 s \\
    \rowcolor{red!12}
    Codex + Qwen     & 30.0\%  & 83.3\%  & \textbf{46.7\%} & 3.53 & 315 s \\
    \hline
  \end{tabular}

  \vspace{2.5mm}
  \statbox{83.3\% / 83.3\%}{DeepSeek 下两框架测试通过率\textbf{完全持平}}
  \hfill
  \statbox{66.7\% vs 46.7\%}{Qwen 下自研\textbf{高出 20 个百分点}}
  \hfill
  \statbox{96.7\% vs 100\%}{首轮编译：自研靠反馈追平最终结果}
\end{frame}

\begin{frame}
  \frametitle{代价：同样的通过率，输入成本差一个量级}

  \centering
  \small
  \setlength{\tabcolsep}{6pt}
  \renewcommand{\arraystretch}{1.2}
  \begin{tabular}{lcccc}
    \hline
    设置 & 输入 tokens & 输出 tokens & 平均耗时 & 相对 Codex \\
    \hline
    自研 + DeepSeek  & 75,058    & 383,312 & 163 s & 输入 \textbf{13\%}、耗时 \textbf{48\%} \\
    自研 + Qwen      & 131,255   & 503,606 & 295 s & 输入 \textbf{8\%}、耗时 \textbf{68\%} \\
    Codex + DeepSeek & 565,110   & 521,147 & 342 s & --- \\
    Codex + Qwen     & 1,404,958 & 394,781 & 315 s & --- \\
    \hline
  \end{tabular}

  \vspace{2.5mm}
  \begin{itemize}
    \item 通过率相同（83.3\%）的前提下，自研输入 token 为 Codex 的 \textbf{1/7}、耗时约 \textbf{1/2}
    \item 差距几乎全在\textbf{输入侧}：Codex 的 agent 循环反复累加系统提示、工具定义与历史
    \item 输出 token 两框架同量级（383k vs 521k），说明“想的量”相当，差的是“看的量”
  \end{itemize}
\end{frame}

\begin{frame}
  \frametitle{解释一：差距落在编译层，不是行为层}

  \begin{columns}[T]
    \begin{column}{0.53\textwidth}
      \centering
      \resizebox{\columnwidth}{!}{%
      \begin{tikzpicture}[x=0.19cm,y=1cm]
        % 图例
        \fill[teal!60!black] (0,0.55) rectangle (2,0.85);
        \node[anchor=west,font=\scriptsize] at (2.2,0.70) {成功};
        \fill[orange!85!black] (9,0.55) rectangle (11,0.85);
        \node[anchor=west,font=\scriptsize] at (11.2,0.70) {编译未通过};
        \fill[red!70!black] (18.5,0.55) rectangle (20.5,0.85);
        \node[anchor=west,font=\scriptsize] at (20.7,0.70) {测试未通过};

        \foreach \i/\name/\a/\b/\c in {%
            0/自研+DeepSeek/25/0/5,
            1/自研+Qwen/20/2/8,
            2/Codex+DeepSeek/25/0/5,
            3/Codex+Qwen/14/5/11}{%
          \node[anchor=east,font=\scriptsize] at (-0.6,{-\i*0.75-0.4}) {\name};
          \fill[teal!60!black]   (0,{-\i*0.75-0.55}) rectangle (\a,{-\i*0.75-0.25});
          \fill[orange!85!black] (\a,{-\i*0.75-0.55}) rectangle ({\a+\b},{-\i*0.75-0.25});
          \fill[red!70!black]    ({\a+\b},{-\i*0.75-0.55})
                                 rectangle ({\a+\b+\c},{-\i*0.75-0.25});
          \node[font=\scriptsize,text=white] at ({\a/2},{-\i*0.75-0.4}) {\a};
          \ifnum\b>0
            \node[font=\scriptsize,text=white] at ({\a+0.5*\b},{-\i*0.75-0.4}) {\b};
          \fi
          \node[font=\scriptsize,text=white] at ({\a+\b+0.5*\c},{-\i*0.75-0.4}) {\c};
        }
        \node[anchor=north west,font=\tiny] at (0,-3.05) {共 30 题，横轴为题目数};
      \end{tikzpicture}}
    \end{column}
    \begin{column}{0.44\textwidth}
      \begin{itemize}
        \item DeepSeek 两种框架\textbf{编译失败均为 0}：
              编译闸门已不是瓶颈，扣分全在行为层
        \item 自研 + Qwen 只剩 2 题编不过，靠反馈把 28 题送进测试
        \item Codex + Qwen 有 \textbf{5 题最终仍编译失败}：
              首轮编译 30\%，5 轮预算常耗在语法上
        \item 只看编译通过的题，行为通过率为 83.3\% / 83.3\% / 71.4\% / 56.0\%：
              Qwen 下两框架的差距从 \textbf{20 个百分点收窄到 15 个百分点}，
              其余差距来自没编译过的题
      \end{itemize}
    \end{column}
  \end{columns}
\end{frame}

\begin{frame}
  \frametitle{解释二：Qwen 下差距从哪来（三组机制性证据）}

  \begin{columns}[T]
    \begin{column}{0.45\textwidth}
      \centering
      \scriptsize
      \setlength{\tabcolsep}{3.5pt}
      \begin{tabular}{@{}lcc@{}}
        \hline
        指标（Qwen） & 自研 & Codex \\
        \hline
        首轮进入编译            & \textbf{30/30} & \textbf{21/30} \\
        \quad 接口预校验拦下    & 0 & 7 \\
        \quad 块解析拦下        & 0 & 2 \\
        每轮编译错误条目        & \textbf{1.17} & \textbf{2.66} \\
        首轮错误数（中位/均值） & 0 / 1.2 & 2 / 3.7 \\
        首轮编译通过            & 23/30 & 9/30 \\
        双方都过的题平均轮次    & \textbf{1.38} & \textbf{2.00} \\
        \hline
      \end{tabular}

      \vspace{2mm}
      \raggedright\scriptsize
      21 道至少一方通过的题：双方都过 13、自研独有 7、Codex 独有 1。
      失败题集高度重合——差距不是“某些题完全不会”，
      而是\textbf{同等能力下 Codex 要更多轮才收敛}（上限 5 轮）。
    \end{column}
    \begin{column}{0.52\textwidth}
      \scriptsize
      \begin{itemize}
        \item \textbf{输出结构不成立}：Codex 侧把 \texttt{REAL}/\texttt{BOOL}
              当 DUT 名、数组别名写成 \texttt、整段没有 POU 块——
              9 题的提交根本没走到编译器
        \item \textbf{错误更像通用语言}：\texttt{C189} 74、\texttt{C9} 68 的样本是
              \texttt{应该用',' 替换')'}、\texttt{应该用表达式代替 '!'}——
              把 ST 写成 C/JS 风格（\texttt{!}、\texttt{if(...)}、\texttt{//}）
        \item \textbf{提示层稀释}：Codex 以通用 agent 身份运行，模型要同时满足
              “我是编码代理”与“输出严格块格式的单个 FB”；Qwen 指令跟随弱，
              冲突时滑向通用写法——这正是 DeepSeek 打平、Qwen 拉开的原因
        \item \textbf{上下文层}：中位输入 token 11,314 vs 1,573（7.2 倍），
              输出预算被压缩（3,354 vs 6,101）
      \end{itemize}
    \end{column}
  \end{columns}

  \vspace{1mm}
  \raggedright\scriptsize
  如实标注：20/30 vs 14/30 的 Fisher 精确检验 $p=0.19$，未达 0.05 显著；
  结论依据上述机制性证据，而非单一通过率（每题仅采样 1 次）。
\end{frame}

\begin{frame}
  \frametitle{剩余失败在哪里：行为语义，编译反馈管不到}

  \begin{columns}[T]
    \begin{column}{0.48\textwidth}
      \begin{itemize}
        \item 状态机题：五工位生产、柔性加工
        \item 定时/边沿题：焊缝去抖、自动门、交通灯
        \item 典型症状：“差一拍”、未覆盖优先级
        \item 2 题在\textbf{全部 4 个设置下同时失败}
              （\texttt{public\_hard\_003}、\texttt{self\_easy\_003}），
              另 3 题在 3 个设置下失败（\texttt{public\_easy\_005}、
              \texttt{public\_hard\_004}、\texttt{public\_medium\_003}）——
              指向用例与规约的判定强度，而非某个模型的能力
      \end{itemize}
    \end{column}
    \begin{column}{0.48\textwidth}
      \begin{alertblock}{判定强度的缺口（如实标注）}
        \begin{itemize}
          \item 驱动只断言 \texttt{expect} 中的输出：
                12 个扫描步无断言
          \item 34 个布尔输入从未置 TRUE，分布在 9 道题
          \item 规约未接入 PLCVerif；每题仅采样 1 次
        \end{itemize}
      \end{alertblock}
      因此上表通过率是\textbf{“现有用例下”}的通过率，严格化后可能下降。
    \end{column}
  \end{columns}
\end{frame}

% ============================================================
\section{总结}

\begin{frame}
  \frametitle{总结}

  \begin{enumerate}
    \item 建成 30 题数据集，逐题 LTL 规约、测试用例（含逐步测试目的）与参考实现；
          参考实现 30/30 通过编译与仿真，形成可复现基线
    \item 自研 harness 在 DeepSeek 下达到与 Codex 相同的测试通过率（83.3\%），
          输入 token 约为其 1/7、耗时约 1/2；在 Qwen 下仍更优（66.7\% vs 46.7\%）
    \item 差距主要落在\textbf{首轮进入编译的比例}（30/30 vs 21/30）、
          每轮错误密度（1.17 vs 2.66）与上下文成本，
          而非“模型想得不够”；剩余失败集中在行为语义
  \end{enumerate}

  \vspace{1mm}
  \begin{block}{如实标注的局限}
    \begin{itemize}
      \item 规约未接入 PLCVerif（选做项，未做）；每题仅采样 1 次，
            Qwen 下两框架差异 Fisher 精确检验 $p=0.19$，未达显著
      \item 12 个扫描步无断言、34 个布尔输入从未置 TRUE，通过率是“现有用例下”的通过率
      \item 下一步：接入 PLCVerif、补强用例覆盖、多平台验证
    \end{itemize}
  \end{block}

  \vspace{0.3em}
  \begin{center}\small
    把“能不能编译、跑得对不对”做成结构化反馈，弱模型也能被推上来；
    但没有对所有模型都更好的框架。
  \end{center}
\end{frame}

\end{document}
